% GENERATED FILE. Edit canonical lessons, then run npm run generate:books.
% canonical-source-hash: fnv1a64:50b1d1f6d6ab2fa6
% canonical-lessons: MR-C261-kurup, MR-C261-thaklela, MR-C261-bhukela, MR-C261-tahanlela, MR-C261-ragit

\chapter{Ugly, Tired, Hungry, Thirsty, Short-tempered}
\label{ch:mr-ugly-tired-hungry-thirsty-short-tempered}
\hlchaptermodality{261}

\begin{chapteropening}
\textbf{By the end of this chapter:} \emph{I can say ugly, tired, hungry, thirsty and short-tempered.}

Complete the last lesson of chapter 261: I can say ugly, tired, hungry, thirsty and short-tempered.
\end{chapteropening}

\section[\texorpdfstring{\mr{कुरूप} (kurūp)}{कुरूप (kurūp)}]{\mr{कुरूप} (kurūp) — ugly}

\label{lesson:MR-C261-kurup}



\begin{quote}
Before the new one: say the Marathi for unmarried, then the Marathi for married.
\end{quote}

\subsection*{You'll want to know: \mr{कुरूप}}
\textbf{\mr{कुरूप}} — \emph{kurūp} — \textquotedblleft{}ugly\textquotedblright{}.

\begin{grammarlens}[title={What you've built}]
One more word for this chapter.
\end{grammarlens}

\subsection*{Your turn}
\begin{itemize}
\raggedright
  \item \emph{Say it:} \emph{kurūp}
  \item \emph{Say it:} \emph{kurūp}, once more
  \item \emph{Say it:} say \emph{kurūp}
  \item \emph{Recall:} say the Marathi for unmarried, then the Marathi for married, then say \emph{kurūp} again
\end{itemize}

\begin{tcolorbox}[breakable,colback=teal!4,colframe=teal!35!black,title={Before you move on}]
What does \mr{कुरूप} mean? (\textquotedblleft{}Ugly\textquotedblright{}.) Say it once more.
\end{tcolorbox}
\section[\texorpdfstring{\mr{थकलेला} (thaklelā)}{थकलेला (thaklelā)}]{\mr{थकलेला} (thaklelā) — tired}

\label{lesson:MR-C261-thaklela}



\begin{quote}
Before the new one: say the Marathi for married, then the Marathi for ugly.
\end{quote}

\subsection*{You'll want to know: \mr{थकलेला}}
\textbf{\mr{थकलेला}} — \emph{thaklelā} — \textquotedblleft{}tired\textquotedblright{}.

\begin{grammarlens}[title={What you've built}]
One more word for this chapter.
\end{grammarlens}

\subsection*{Your turn}
\begin{itemize}
\raggedright
  \item \emph{Say it:} \emph{thaklelā}
  \item \emph{Say it:} \emph{thaklelā}, once more
  \item \emph{Say it:} say \emph{thaklelā}
  \item \emph{Recall:} say the Marathi for married, then the Marathi for ugly, then say \emph{thaklelā} again
\end{itemize}

\begin{tcolorbox}[breakable,colback=teal!4,colframe=teal!35!black,title={Before you move on}]
What does \mr{थकलेला} mean? (\textquotedblleft{}Tired\textquotedblright{}.) Say it once more.
\end{tcolorbox}
\section[\texorpdfstring{\mr{भुकेला} (bhukelā)}{भुकेला (bhukelā)}]{\mr{भुकेला} (bhukelā) — hungry}

\label{lesson:MR-C261-bhukela}



\begin{quote}
Before the new one: say the Marathi for ugly, then the Marathi for tired.
\end{quote}

\subsection*{You'll want to know: \mr{भुकेला}}
\textbf{\mr{भुकेला}} — \emph{bhukelā} — \textquotedblleft{}hungry\textquotedblright{}.

\begin{grammarlens}[title={What you've built}]
One more word for this chapter.
\end{grammarlens}

\subsection*{Your turn}
\begin{itemize}
\raggedright
  \item \emph{Say it:} \emph{bhukelā}
  \item \emph{Say it:} \emph{bhukelā}, once more
  \item \emph{Say it:} say \emph{bhukelā}
  \item \emph{Recall:} say the Marathi for ugly, then the Marathi for tired, then say \emph{bhukelā} again
\end{itemize}

\begin{tcolorbox}[breakable,colback=teal!4,colframe=teal!35!black,title={Before you move on}]
What does \mr{भुकेला} mean? (\textquotedblleft{}Hungry\textquotedblright{}.) Say it once more.
\end{tcolorbox}
\section[\texorpdfstring{\mr{तहानलेला} (tahānlelā)}{तहानलेला (tahānlelā)}]{\mr{तहानलेला} (tahānlelā) — thirsty}

\label{lesson:MR-C261-tahanlela}



\begin{quote}
Before the new one: say the Marathi for tired, then the Marathi for hungry.
\end{quote}

\subsection*{You'll want to know: \mr{तहानलेला}}
\textbf{\mr{तहानलेला}} — \emph{tahānlelā} — \textquotedblleft{}thirsty\textquotedblright{}.

\begin{grammarlens}[title={What you've built}]
One more word for this chapter.
\end{grammarlens}

\subsection*{Your turn}
\begin{itemize}
\raggedright
  \item \emph{Say it:} \emph{tahānlelā}
  \item \emph{Say it:} \emph{tahānlelā}, once more
  \item \emph{Say it:} say \emph{tahānlelā}
  \item \emph{Recall:} say the Marathi for tired, then the Marathi for hungry, then say \emph{tahānlelā} again
\end{itemize}

\begin{tcolorbox}[breakable,colback=teal!4,colframe=teal!35!black,title={Before you move on}]
What does \mr{तहानलेला} mean? (\textquotedblleft{}Thirsty\textquotedblright{}.) Say it once more.
\end{tcolorbox}
\section[\texorpdfstring{\mr{रागीट} (rāgīṭ)}{रागीट (rāgīṭ)}]{\mr{रागीट} (rāgīṭ) — short-tempered}

\label{lesson:MR-C261-ragit}



\begin{quote}
Before the new one: say the Marathi for hungry, then the Marathi for thirsty.
\end{quote}

\subsection*{You'll want to know: \mr{रागीट}}
\textbf{\mr{रागीट}} — \emph{rāgīṭ} — \textquotedblleft{}short-tempered\textquotedblright{}.

\begin{grammarlens}[title={What you've built}]
One more word for this chapter.
\end{grammarlens}

\subsection*{Your turn}
\begin{itemize}
\raggedright
  \item \emph{Say it:} \emph{rāgīṭ}
  \item \emph{Say it:} \emph{rāgīṭ}, once more
  \item \emph{Say it:} say \emph{rāgīṭ}
  \item \emph{Recall:} say the Marathi for hungry, then the Marathi for thirsty, then say \emph{rāgīṭ} again
\end{itemize}

\begin{tcolorbox}[breakable,colback=teal!4,colframe=teal!35!black,title={Before you move on}]
What does \mr{रागीट} mean? (\textquotedblleft{}Short-tempered\textquotedblright{}.) Say it once more.
\end{tcolorbox}
